\documentclass[12pt]{article}
\usepackage[T1]{fontenc}
\usepackage[utf8]{inputenc}
\usepackage{lmodern}
\usepackage{textcomp}
\usepackage{enumitem}
\usepackage{geometry}
\geometry{margin=1in}
\begin{document}
\begin{center}
Answer Key - Political Science - Undergraduate Level Assessment 3 \
\small (Answers are ordered exactly as in the question paper.)
\end{center}

\section*{Multiple Choice}
\begin{enumerate}[label=\arabic*.]
\item \textbf{Answer:} D
\\ \textit{Options:} A) Congress asserted its primacy in foreign policy; B) US foreign policy became substantially decentralized; C) The Presidency welcomed the influence of Congress; D) None of the above
\\ \textit{Note:} The correct answer is "None of the above" because the relationship between President George H.W. Bush and Congress was characterized by cooperation during the early years followed by increased partisanship, while Bill Clinton faced significant legislative challenges that reflected a contentious dynamic, thus making all provided options inaccurate.
\item \textbf{Answer:} D
\\ \textit{Options:} A) Plessy v. Ferguson decision; B) Missouri Compromise; C) Jim Crow laws; D) Fourteenth Amendment
\\ \textit{Note:} The Dred Scott decision, which denied citizenship and constitutional rights to African Americans, was effectively overturned by the Fourteenth Amendment, which guaranteed citizenship and equal protection under the law to all individuals born or naturalized in the United States.
\item \textbf{Answer:} D
\\ \textit{Options:} A) The external environment of the international system; B) Domestic environment, government and bureaucracy; C) The President's individual personality.; D) All of the above
\\ \textit{Note:} The correct answer "All of the above" is valid because Rosenau emphasizes the interplay of multiple factors, including domestic politics, international pressures, and historical context, in shaping the foreign policy behavior of the United States.
\item \textbf{Answer:} A
\\ \textit{Options:} A) Coining money; B) Authorizing constitutional amendments; C) Having representation in Congress; D) Appealing to the president to adjudicate disputes
\\ \textit{Note:} The Articles of Confederation granted states the power to coin money independently, while the Constitution centralized this authority under the federal government, prohibiting states from issuing their own currency.
\item \textbf{Answer:} C
\\ \textit{Options:} A) release detainees from prison following unlawful arrest; B) maintain the navy; C) increase the power of the national government; D) veto legislative bills of attainder
\\ \textit{Note:} The Commerce Clause has been interpreted by the Supreme Court to grant Congress the authority to regulate interstate commerce, thereby expanding the scope of federal power and enabling the national government to enact laws affecting economic activities across state lines.
\end{enumerate}

\section*{True / False}
\begin{enumerate}[label=\arabic*.]
\setcounter{enumi}{5}
\item \textbf{Answer:} True
\\ \textit{Options:} A) True; B) False
\\ \textit{Note:} Voters aged 18-29 are generally the demographic least likely to participate in elections due to factors such as lower political engagement, varying levels of education, and a lack of established voting habits compared to older age groups.
\item \textbf{Answer:} True
\\ \textit{Options:} A) True; B) False
\\ \textit{Note:} The statement is true because Critical Security Studies challenges traditional notions of security that prioritize state-centric views and instead emphasizes a variety of referent objects, including individuals, communities, and broader social constructs, thus allowing for diverse interpretations of what security entails.
\item \textbf{Answer:} False
\\ \textit{Options:} A) True; B) False
\\ \textit{Note:} The statement is false because the U.S. Constitution grants states the authority to determine voter eligibility, although Congress can regulate certain aspects of the electoral process under its constitutional powers.
\item \textbf{Answer:} True
\\ \textit{Options:} A) True; B) False
\\ \textit{Note:} The death penalty is considered by some to be cruel and unusual punishment because it raises ethical concerns about the potential for wrongful executions, disproportionate application based on race or socioeconomic status, and the inhumane methods of execution that can cause suffering.
\item \textbf{Answer:} True
\\ \textit{Options:} A) True; B) False
\\ \textit{Note:} The First Amendment explicitly guarantees the right of American citizens to petition the government, ensuring that individuals can seek remedies for grievances and communicate their concerns to governmental authorities.
\end{enumerate}

\section*{Long Form Response}
\begin{enumerate}[label=\arabic*.]
\setcounter{enumi}{10}
\item \textbf{Answer:} Donald Trump's 2016 campaign rhetoric effectively framed globalization as harmful to 'ordinary Americans,' portraying it as a force that stripped away jobs and economic stability from the working class. He argued that liberal elites were the architects of globalization, benefiting from international trade agreements and economic policies while disregarding the negative impacts on everyday citizens. This narrative resonated with many voters who felt left behind by economic changes, as Trump emphasized the loss of manufacturing jobs and the decline of local economies. By positioning himself as a champion for the 'forgotten man' against the interests of the elite, Trump tapped into a widespread sentiment of disenfranchisement, ultimately galvanizing support for his campaign. This framing of globalization underscored a broader critique of political establishments and their perceived disconnect from the realities faced by ordinary Americans.
\\ \textit{Note:} correct because it accurately captures how Trump's rhetoric positioned globalization as a threat to the economic security of ordinary Americans, attributing the responsibility for these adverse effects to liberal elites who were perceived as profiting from policies that neglected the working class.
\item \textbf{Answer:} The relationship between the state-centric approach and human security reveals significant limitations in prioritizing state security over individual well-being. While state-centric frameworks focus on the protection and stability of the state itself, human security shifts the focus to the safety and dignity of individuals and communities. This prioritization of state interests can lead to neglect of pressing human rights issues, social injustices, and humanitarian crises that affect populations. Consequently, by emphasizing state security, policies may overlook essential aspects of human security, such as economic stability, health, and personal safety, ultimately compromising the well- being of individuals. Thus, adopting a human security framework underscores the importance of placing individual needs at the forefront of security discussions, highlighting the inadequacies of a purely state-centric perspective.
\\ \textit{Note:} The answer correctly identifies the fundamental tension between a state-centric security approach, which prioritizes the state's stability, and the human security perspective, which emphasizes the importance of individual well-being and rights, thereby illustrating the potential neglect of critical humanitarian issues when state interests take precedence.
\end{enumerate}

\vfill
\noindent\textit{End of Answer Key}
\end{document}